% SOURCE_AUTHORITY: sga5_fr_workpass.tex, lines 8469--8520 (LF-normalized unit).
% SOURCE_SHA256: 791F4EFFC5E02832D5D77ED03518C8156D6F07E4C8238B03545DB93D883FBB28
\begin{statement}{Proposição 1.3.2}
Se \(F\) e \(G\) forem dois \(\mathbb Z_\ell\)-feixes construtíveis, o produto tensorial \(F\otimes G\) é um \(\mathbb Z_\ell\)-feixe construtível.
\end{statement}

\begin{prooftext}[Demonstração]
A \(\ell\)-adicidade é clara. Por outro lado, segue-se imediatamente da definição dos feixes construtíveis (SGA 4 IX 2.3) que o produto tensorial de dois feixes construtíveis é construtível.
\end{prooftext}

Por restrição define-se, portanto, um bifuntor aditivo exato à direita, também chamado produto tensorial:
\[
\otimes:\mathbb Z_\ell\text{-}\fc(X)\times \mathbb Z_\ell\text{-}\fc(X)
\longrightarrow \mathbb Z_\ell\text{-}\fc(X).
\]

\subsubsection*{1.3.3. Feixe de homomorfismos}

Sejam \(F=(F_n)_{n\in\mathbb N}\) e \(G=(G_n)_{n\in\mathbb N}\) dois \(\mathbb Z_\ell\)-feixes construtíveis. Supõe-se que \(F\) é localmente livre, entendendo por isso que, para todo inteiro \(n\), o feixe \(F_n\) é localmente livre sobre o feixe constante
\[
(\mathbb Z/\ell^{\,n+1}\mathbb Z)_X.
\]
Para todo inteiro \(n\), a sequência exata estrutural
\[
G_{n+1}\xrightarrow{\ell^{\,n+1}}G_{n+1}\longrightarrow G_n\longrightarrow 0
\]
fornece, como \(F_{n+1}\) é localmente livre sobre \((\mathbb Z/\ell^{\,n+1}\mathbb Z)_X\), uma sucessão exata
\[
\cHom(F_{n+1},G_{n+1})
 \xrightarrow{\ell^{\,n+1}}
\cHom(F_{n+1},G_{n+1})
 \longrightarrow
\cHom(F_{n+1},G_n)
 \longrightarrow 0.
\tag{1}
\]
Como \(\ell^{\,n+1}G_n=0\), a sequência exata estrutural
\[
F_{n+1}\xrightarrow{\ell^{\,n+1}}F_{n+1}\longrightarrow F_n\longrightarrow 0
\]
define claramente um isomorfismo
\[
\cHom(F_{n+1},G_n)\xleftarrow{\sim}\cHom(F_n,G_n).
\tag{2}
\]
Combinando (1) e (2), obtém-se, para todo inteiro \(n\), um morfismo
\[
\cHom(F_{n+1},G_{n+1})\longrightarrow \cHom(F_n,G_n).
\]
Por isso, um sistema projetivo
\[
\cHom(F,G)=\bigl(\cHom(F_n,G_n)\bigr)_{n\in\mathbb N},
\]
que é claramente um \(\mathbb Z_\ell\)-feixe.
